% ======================================================================
% sections/introduction.tex
%
% Introduction for the full paper:
%   Triple Humanoid 24/7 Adverse Event Oncology Trial Response Team:
%   4-Site Rotation (v0.7.0)
%
% Five subsections inside the single \section{Introduction} block:
%   1) Background and Setting
%   2) Thesis Statement and Scope
%   3) The On-Premises Repository LLM Pattern
%   4) The 4-Site Continental Rotation
%   5) Roadmap of the Paper
% ======================================================================

\begin{figure}[H]
    \hspace{-0.95cm}
   \includegraphics[width=1.1\linewidth, trim={0 125 0 0}, clip]{Images/fig-01-network-layout.png}
    % \vspace{-0.25cm}
    \caption{Teal=per-site LLM, Gray=broadcast bus, Blue=lead H2, Aqua=assist H2, Green=reserve H2.}
    \label{Overview1}
\end{figure}


\begin{sectionaccent}
\subsection{Background and Setting}
\label{sec:intro-background}

A prior 2025 author paper demonstrated that large language models are
effective at iterative learning against the FDA Adverse Event
Reporting System (FAERS) when 16 proprietary and open source models
are placed in head to head competition over the same code generation
task \cite{kawchak_2025_18029100}. That paper established the
plausibility of an autonomous LLM pipeline against FDA adverse event
data. The current paper does not redo that demonstration. Instead, the current
paper extends it from a single agent FAERS task to a coordinated
3 robot camarade swarm 24/7 adverse event response team running at
4 trial sites at once.

\subsection{Thesis Statement and Scope}
\label{sec:intro-thesis}

The thesis sentence is preserved verbatim from the project
instructions tree \cite{kawchak_2026_humanoid_instructions}:
on-premises repository based LLMs provide commands to humanoid robots
based on real-time sensor data and controlled via x, y, z coordinates
to administer synergistic treatment regarding patient adverse events;
this workflow minimizes single robot error potential. The thesis
applies to a defined scope shown in \autoref{Overview1}: 4 trial sites under the network
identifier PAT-NET-001, 3 Unitree H2 EDU humanoid robots per site, 1
on-premises Claude Opus 4.7 1M Max language model per site, a 168 hour
monitoring window, and a 32 iteration sweep across the
language model temperature, the adverse event arrival jitter, the
response time variance, and the camarade aggressiveness axes
\cite{unitree_h2_2026, nvidia_jetson_thor_2026,
anthropic_claude_opus_47_2026}.

The scope statement matters. The paper does not claim that a live
Unitree H2 EDU body executed any broadcast command. It does not claim
that real adverse event records were submitted to the FDA RTCT pilot
intake API. It does claim that the full software pipeline, the
schemas, the kinematics constraints, the audit chain, and the safety
envelope are present, tested, and ready for the governance step that
the technology layer cannot perform on its own.

\subsection{The On-Premises Repository LLM Pattern}
\label{sec:intro-onprem}

The pattern centers on one Claude Opus 4.7 1M Max instance hosted on
a site local server. The site model reads a frozen system prompt, a
defined input schema for sensor frames, and a defined output schema
for a structured language model decision payload. At every 1 Hz tick
the model emits exactly 3 sub commands, one for each robot at the
site. The broadcaster then publishes those 3 sub commands as a single
atomic message that all 3 robots subscribe to in parallel through a
publish subscribe channel. The on-premises placement keeps inference
latency near a 200 ms median over a local network and keeps every
inference traceable to a known repository version
\cite{anthropic_claude_code_2026, anthropic_claude_opus_47_2026}.
Robots talk to each other through two channels at once: a 60 GHz
ultra wide band peer mesh for short range physical handoff
coordination, paired with an infrared beacon line of sight channel
for visual cross check, and a shared on-prem compute fabric for
intellectual exchange of sensor evidence and role decisions
\cite{kawchak_2026_humanoid_instructions}. The cross site bus carries
only de identified hourly summaries across the 4 sites; no robot
transits, and no protected health information.

\subsection{The 4-Site Continental Rotation}
\label{sec:intro-rotation}

The 4 site continental network spans San Francisco (site SF-01), San
Diego (site SD-01), Boston (site BO-01), and Atlanta (site AT-01).
Each site holds a 3 robot Unitree H2 EDU swarm with the lead, assist,
and reserve roles, with roles rotated every 1 Hz tick seen in \autoref{3H2s}. Figure
\ref{fig:network-layout} carries the verbatim layout diagram from
the codegen tree. Figure \ref{fig:execution-sequence} carries the
verbatim 4 site execution sequence diagram from the execution tree
and shows the per site dispatch, the 1 Hz language model decisions,
and the audit chain extension at every tick.

\clearpage

\begin{figure}[H]
\vspace{-0.15cm}
\centering
\begin{Verbatim}[fontsize=\small]
                  PAT-NET-001 Continental Adverse Event Response Network
                  =======================================================

                                  Cross Site Coordination Bus
                                  (de-identified summaries)
                       +-------------+--------------+--------------+
                       |             |              |              |             
                       v             v              v              v             
                  +-------+      +-------+      +-------+      +-------+
                  | SF-01 |      | SD-01 |      | BO-01 |      | AT-01 |
                  +-------+      +-------+      +-------+      +-------+
            
                  Per-Site Layout (same pattern at all 4 sites):
            
                  +--------------------------------------+
                  | Site <ID>                            |
                  | Claude Opus 4.7 1M on-prem           |
                  | Broadcast bus at 1 Hz                |
                  |                                      |
                  |   +------+  +------+   +------+      |
                  |   | H2-1 |  | H2-2 |   | H2-3 |      |
                  |   | Lead |  |Assist|   |Reserve      |
                  |   +------+  +------+   +------+      |
                  +--------------------------------------+

                  Robot ID assignments by site:
                    SF-01: H2-A (Lead), H2-B (Assist), H2-C (Reserve)
                    SD-01: H2-D (Lead), H2-E (Assist), H2-F (Reserve)
                    BO-01: H2-G (Lead), H2-H (Assist), H2-I (Reserve)
                    AT-01: H2-J (Lead), H2-K (Assist), H2-L (Reserve)

                  Legend:
                    Roles shown are typical at AE start.
                    Roles rotate at every 1 Hz tick.
                    Inter robot: 60 GHz UWB peer mesh plus IR beacon line of sight.
                    LLM to robot: 1 Hz broadcast publish subscribe.
                    Cross site: de-identified hourly summaries only.
                    No PHI. No robot transit between sites.
                    Platform: Unitree H2 EDU.
                    Compute upgradeable to Jetson AGX Thor 2070 TOPS.
\end{Verbatim}
\caption{PAT-NET-001 layout with one Claude Code and 3 Unitree H2 EDU robots per site.}
\label{fig:network-layout}
\end{figure}


\begin{figure}[H]
    % \hspace{-1.75cm}
   \includegraphics[width=1\linewidth, trim={0 0 0 0}, clip]{Images/07-10-decision-tree-flowchart.png}
    % \vspace{-0.25cm}
    \caption{Rotation decision tree regarding three H2 robots clinical trial availability.}
    \label{3H2s}
\end{figure}



\begin{figure}[H]
\vspace{-0.15cm}
\centering
\begin{Verbatim}[fontsize=\small]
                     PAT-NET-001 4 Site Execution Sequence (Per Tick)
                     ================================================

       wall clock    | SF-01          | SD-01          | BO-01          | AT-01
       --------------+----------------+----------------+----------------+--------
       T+0    site   | dispatch init  | dispatch init  | dispatch init  | dispat
                 | 3 robots: ABC  | 3 robots: DEF  | 3 robots: GHI  | 3: JKL
                 |                |                |                |
       T+0.0  tick 0 | LLMPlanner 1Hz | LLMPlanner 1Hz | LLMPlanner 1Hz | LLMPlnr
                 | -> Broadcaster | -> Broadcaster | -> Broadcaster | -> Bcst
                 | -> Dispatcher  | -> Dispatcher  | -> Dispatcher  | -> Disp
                 | 3 acks 50 ms   | 3 acks 50 ms   | 3 acks 50 ms   | 3 acks
                 |                |                |                |
       T+1.0  tick 1 | sensor refresh | sensor refresh | sensor refresh | sensor
                 | role rebalance | role rebalance | role rebalance | rebal.
                 |                |                |                |
       T+59   tick59 | audit head=    | audit head=    | audit head=    | head=
                 | 6274ffb43...   | 93e5b2a36...   | f26259454...   | 9ce574
                 |                |                |                |
       T+60   summ   | 60 decisions   | 60 decisions   | 60 decisions   | 60 dec
                 | persisted      | persisted      | persisted      | persis
                            
                                             |
                                             v
                           +------------------------------------+
                           | Aggregate                          |
                           | week_runner: 3600 ticks per site   |
                           | iterate.py: 32 iterations          |
                           | runner (Rust release): 32 sweeps   |
                           | average p50 response time 67.5 s   |
                           | camaraderie pass rate 0.985        |
                           +------------------------------------+
\end{Verbatim}
\caption{4 site execution sequence at the 1 Hz language model tick boundary with audit chain extension and the 60 second per site summary.}
\label{fig:execution-sequence}
\end{figure}

The continental rotation is not a geographic load balance. Each site
holds local sensor frames, local patient state, and local physician
state. The cross site bus only carries de identified hourly
summaries. This separation honors patient privacy while still allowing
the 4 sites to share lessons learned across the network through the
shared on prem compute fabric.

\subsection{Roadmap of the Paper}
\label{sec:intro-roadmap}

The methods section walks 9 subsections that anchor every claim to a
specific source file in the read only RESEARCH trees. The results
section contains 7 subsections that present the per site decision
streams, the 32 iteration sweep aggregate, the comparison reports,
the ASCII visuals, the converted output tables, the figures, and the
headline metrics. The discussion section includes 6 subsections that
reconcile each clause of the thesis with the work and pose the
practical insight on real life feasibility. The limitations and future
work section names the simulator versus real robot gap, the per site
data identity bound, the excluded extra hours dataset, the missing
repository update instruction directory, and the roadmap items that
the project hands off to the agencies, sponsors, and sites. The
conclusions section restates the thesis verbatim and presents the
work as an industry checklist already executed.

\subsubsection*{Real life feasibility insight.} As of May 2026 the
Unitree H2 EDU is shipping at the spec sheet level named in the
instructions tree (39 DOF, 1.8 m, 70 kg, 10 kg payload per arm, IP65
outdoor rated) \cite{unitree_h2_2026}. Pairing it with one on
premises Claude Opus 4.7 1M Max instance per site is feasible on
conventional high end x86 or Apple Silicon servers
\cite{anthropic_claude_opus_47_2026}. The 1 Hz broadcast cadence is
well within current network and inference latency budgets. What
remains is the IRB approval, sponsor agreement, and FDA RTCT
pilot enrollment. The paper treats those as part of the forward
roadmap, but not solved.

\sectiontable{Introduction planning matrix}{%
Context & I1 & 1 & FDA RTCT, CTCAE & State setting & 1 to 2 paragraphs & FAERS cite & Done \\
Thesis & I2 & 2 & instructions & Verbatim quote & 1 paragraph & Same as README & Done \\
Pattern & I3 & 3 & llm planner & On-prem LLM & 2 paragraphs & 1 Hz cadence & Done \\
Network & I4 & 4 & net layout & ASCII Figure 1, 2 & 2 figures & 4 sites visible & Done \\
Roadmap & I5 & 5 & main.tex & Section walk & 1 paragraph & Reader path & Done \\
}


\end{sectionaccent}
